% !TEX TS-program = pdflatex
% !TeX program = pdflatex
% !TEX encoding = UTF-8
% !TEX spellcheck = en_US

\documentclass[xcolor=table]{beamer}


%\usepackage{fullpage}
%\usepackage[left=2.8cm,right=2.2cm,top=2 cm,bottom=2 cm]{geometry}
\setbeamersize{text margin left=10pt,text margin right=10pt}



%\usepackage[T1]{fontenc}
\usepackage{textcomp}
\usepackage[utf8]{inputenc}
%\usepackage[french,english]{babel}
\usepackage{arabtex}
\usepackage{txfonts}
\usepackage{tipa}
\usepackage[]{graphicx}
\usepackage{multirow}
\usepackage{hyperref}
\usepackage{colortbl}
\usepackage{tabularray}
\usepackage{listingsutf8}
\usepackage{wrapfig}
\usepackage{multicol}
\usepackage[export]{adjustbox} %for images in table, also for frame
\usepackage[many]{tcolorbox}
\usepackage{wasysym}
\usepackage[lined]{algorithm2e}
\usepackage{alltt} %verbatim with commands
\usepackage{longtable}
\usepackage{tabu}

\usepackage{CJKutf8}

\usepackage{amsmath,amssymb, amsfonts} 

%\usepackage{newtxtext,newtxmath}


%\usepackage{acolor}


\definecolor{lightblue}{HTML}{D0D2FF}
\definecolor{lightyellow}{HTML}{FFFFAA}
\definecolor{darkblue}{HTML}{0000BB}
\definecolor{olivegreen}{HTML}{006600}
\definecolor{darkgreen}{HTML}{008B45} %009B55
\definecolor{violet}{HTML}{6600CC}
\definecolor{deeppink}{HTML}{FF1493}
\definecolor{orangey}{HTML}{FFBB00}


\makeatletter
\def\beamer@calltheme#1#2#3{%
	\def\beamer@themelist{#2}
	\@for\beamer@themename:=\beamer@themelist\do
	{\IfFileExists{\beamer@themelocation/#3\beamer@themename.sty}{%
	\usepackage[{#1}]{\beamer@themelocation/#3\beamer@themename}}{\usepackage[{#1}]{#3\beamer@themename}}
		
	}%
}

\def\usefolder#1{
	\def\beamer@themelocation{#1}
}
\def\beamer@themelocation{}
\makeatother

\usefolder{../../extra/beamer-themes}

\usetheme{Karim} % Antibes Boadilla Warsaw Karim

\beamertemplatenavigationsymbolsempty

\let\oldcite\cite
\renewcommand{\cite}[1]{{\bfseries\color{blue}\oldcite{#1}}}

\tcbuselibrary{listings}


%\renewcommand{\baselinestretch}{1.5}

\def\supit#1{\raisebox{0.8ex}{\small\it #1}\hspace{0.05em}}

\AtBeginSection{%
	\begin{frame}
		\sectionpage
	\end{frame}
}

\newcommand{\rottext}[2]{%
	\rotatebox{90}{%
	\begin{minipage}{#1}%
		\raggedleft#2%
	\end{minipage}%
	}%
}


\institute{ %
Laboratoire de la Communication dans les Systèmes Informatiques (LCSI)
\par
École  nationale Supérieure d'Informatique (ESI), Algiers, Algeria
}
\author[Abdelkrime Aries (2023/2024)] %
{Abdelkrime Aries}
%\titlegraphic{\includegraphics[height=1cm]{../img/esi-logo.png}%\hspace*{4.75cm}~


\date{Academic year: 2023/2024} %\today

\titlegraphic{%
	\includegraphics[height=1cm]{../../img/misc/esi-logo.png}%
	\hspace{2cm}%
	\includegraphics[height=1cm]{../../img/misc/lcsi-logo.png}%
	\hspace{2cm}%
	\includegraphics[height=1.5cm]{../../img/misc/esi.nlp-logo.png}%
}

%\setbeamertemplate{headline}{}

\newcommand{\kurl}[1]{{\scriptsize\bfseries\color{orangey}\url{#1}}}

\newcommand{\keyword}[1]{\textcolor{red}{\bfseries\itshape #1}}
\newcommand{\expword}[1]{\textcolor{olivegreen}{#1}}
\newcommand{\optword}[1]{\textcolor{violet}{\bfseries #1}}

\makeatletter
\newcommand\mysphere{%
	\parbox[t]{10pt}{\raisebox{0.2pt}{\beamer@usesphere{item projected}{bigsphere}}}}
\makeatother

%\let\oldtabular\tabular
%\let\endoldtabular\endtabular
%\renewenvironment{tabular}{\rowcolors{2}{white}{lightblue}\oldtabular\rowcolor{blue}}{\endoldtabular}


%\NoAutoSpacing %french autospacing after ":"

\def\graphpath{}

\newcommand{\changegraphpath}[1]{\def\graphpath{#1}}


\newcommand{\vgraphpage}[2][.84\textheight]{%
%	\begin{center}%
		\includegraphics[height=#1]{\graphpath #2}%
%	\end{center}%
}

\newcommand{\hgraphpage}[2][\textwidth]{%
%	\begin{center}%
		\includegraphics[width=#1]{\graphpath #2}%
%	\end{center}%
}

\newcommand{\graphpage}[2][]{%
	\includegraphics[#1]{\graphpath #2}%
}

\bibliographystyle{apalike}

\newcommand{\insertbibliography}[2]{
	\appendix
	\section*{Bibliography}
	\nocite{#2}
%	\makeatletter % to change template
%	\setbeamertemplate{headline}[default] % not mandatory, but I though it was better to set it blank
%	\def\beamer@entrycode{\vspace*{-\headheight}} % here is the part we are interested in :)
%	\makeatother
	\begin{multicols*}{2}[\frametitle{\insertsection} \usebeamertemplate*{frametitle}]%\usebeamertemplate*{frametitle}\frametitle{Références}
		\tiny
		\bibliography{#1}
	\end{multicols*}
}

\definecolor{my-grey}{RGB}{233, 233, 233}

\newcommand{\insertlicence}{
	\begin{frame}[plain]
	\frametitle{License: CC-BY 4.0}

	\begin{tcolorbox}[colback=cyan,
		colframe=cyan,  
		arc=0pt,outer arc=0pt,
		valign=top, 
		halign=center,
		width=\textwidth]
		
		\includegraphics[width=.5cm]{../../img/licence/cc_icon_white_x2.png}
		\includegraphics[width=.5cm]{../../img/licence/attribution_icon_white_x2.png}
		
		\color{white}
		\bfseries Attribution 4.0 International (CC BY 4.0) \\
		\tiny \url{https://creativecommons.org/licenses/by/4.0/deed.en}
		
	\end{tcolorbox}\vspace{-.5cm}
	\begin{tcolorbox}[colback=my-grey,
		colframe=my-grey,  
		center, arc=0pt,outer arc=0pt,
		valign=top, 
		halign=left,
		width=\textwidth]
		
		\tiny
		
		\begin{center}
			\bfseries\Large
			You are free to:
		\end{center}
		
		\begin{minipage}{0.83\textwidth}
			\begin{itemize}
				\item[] \textbf{Share} — copy and redistribute the material in any medium or format
				\item[] \textbf{Adapt} — remix, transform, and build upon the material
				for any purpose, even commercially.
			\end{itemize}
		
			\vspace{12pt}
		
			\textit{The licensor cannot revoke these freedoms as long as you follow the license terms.}
		\end{minipage}
		\begin{minipage}{0.15\textwidth}
			\includegraphics[width=\textwidth]{../../img/licence/FreeCulturalWorks_seal_x2.jpg}
		\end{minipage}
	
		
		\begin{center}
			\bfseries\Large
			Under the following terms:
		\end{center}
		
		\begin{itemize}
			\item[] \textbf{Attribution} — You must give appropriate credit, provide a link to the license, and indicate if changes were made. You may do so in any reasonable manner, but not in any way that suggests the licensor endorses you or your use.
			\item[] \textbf{No additional restrictions} — You may not apply legal terms or technological measures that legally restrict others from doing anything the license permits.
		\end{itemize}
		
	\end{tcolorbox}
	
%	\begin{center}
%		\bfseries Attribution 4.0 International (CC BY 4.0)
%		\url{https://creativecommons.org/licenses/by/4.0/deed.fr}
%	\end{center}

%	\tiny
%
%	Vous êtes autorisé à : 
%	\begin{itemize}
%		\item \textbf{Partager} — copier, distribuer et communiquer le matériel par tous moyens et sous tous formats
%		\item \textbf{Adapter} — remixer, transformer et créer à partir du matériel
%	\end{itemize}
%	
%	Selon les conditions suivantes : 
%	\begin{itemize}
%		\item \textbf{Attribution} — Vous devez créditer l'Œuvre, intégrer un lien vers la licence et indiquer si des modifications ont été effectuées à l'Oeuvre. Vous devez indiquer ces informations par tous les moyens raisonnables, sans toutefois suggérer que l'Offrant vous soutient ou soutient la façon dont vous avez utilisé son Oeuvre.
%		\item \textbf{Pas d'Utilisation Commerciale} — Vous n'êtes pas autorisé à faire un usage commercial de cette Oeuvre, tout ou partie du matériel la composant. 
%		\item \textbf{Pas de restrictions complémentaires} — Vous n'êtes pas autorisé à appliquer des conditions légales ou des mesures techniques qui restreindraient légalement autrui à utiliser l'Oeuvre dans les conditions décrites par la licence.
%	\end{itemize}

	\end{frame}
}

\settowidth{\leftmargini}{\usebeamertemplate{itemize item}}
\addtolength{\leftmargini}{\labelsep}

\AtBeginDocument{
	\newcolumntype{L}[2]{>{\vbox to #2\bgroup\vfill\flushleft}p{#1}<{\egroup}} 
	
	\begin{frame}[plain]
		\maketitle
	\end{frame}

	\insertlicence
}


% needs etoolbox; to break links after -
\appto\UrlBreaks{\do\-}


\makeatletter
\newcommand{\xRightarrow}[2][]{\ext@arrow 0359\Rightarrowfill@{#1}{#2}}
\makeatother


\usefonttheme{structurebold}
%\usefonttheme{professionalfonts}


% figure caption
\setlength\abovecaptionskip{2pt}

% equations
\AtBeginDocument{
\abovedisplayskip=0pt
\abovedisplayshortskip=0pt
\belowdisplayskip=0pt
\belowdisplayshortskip=0pt
}


\title[ESI - NLP: 10- Discourse coherence]%
{Natural Language Processing\\Chapter 10\\Discourse coherence} 

\changegraphpath{../img/coherence/}

\begin{document}
	
\begin{frame}
\frametitle{Natural Language Processing}
\framesubtitle{Discourse coherence: Introduction}

\begin{exampleblock}{Example of a paragraph}
	\begin{center}
		\Large\bfseries
		Winter is one of the year's four seasons.
		The report was well presented.
		Rain is one of this season's qualities.
		I read a report on this season.
	\end{center}
\end{exampleblock}

\begin{itemize}
	\item Can you figure out what this is about?
	\item If there is a problem, what is it?
	\item Can this text be improved?
\end{itemize}

\end{frame}

\begin{frame}
	\frametitle{Natural Language Processing}
	\framesubtitle{Discourse coherence}
	
	\begin{itemize}
		\item How to detect discourse coherence?
		\item \optword{Relations-based coherence}: a discourse sentences are bound together by relations. 
		\begin{itemize}
			\item Rhetorical Structure Theory (RST) 
			\item Penn Discourse TreeBank (PDTB)
		\end{itemize}
		\item \optword{Entity-based coherence}: a discourse sentences form a single entity; they discuss the same subject.
		\begin{itemize}
			\item Centering Theory 
		\end{itemize}
		\item \optword{Lexical cohesion}: close sentences share a few words or meanings.
		\begin{itemize}
			\item Similarity measures
		\end{itemize}
	\end{itemize}
	
\end{frame}

\begin{frame}
\frametitle{Natural Language Processing}
\framesubtitle{Discourse coherence: Some humor}

\begin{center}
	\vgraphpage{humor/humor1.jpg}
\end{center}

\end{frame}

\begin{frame}
\frametitle{Natural Language Processing}
\framesubtitle{Discourse coherence: Plan}

\begin{multicols}{2}
%	\small
\tableofcontents
\end{multicols}
\end{frame}

%===================================================================================
\section{Coherence relations}
%===================================================================================

\subsection{Rhetorical Structure Theory (RST)}

\begin{frame}
\frametitle{Discourse coherence: Coherence relations}
\framesubtitle{Rhetorical Structure Theory (RST)}
	
	\begin{itemize}
		\item Structure  
		\begin{itemize}
			\item \optword{Nucleus (N)} : independent unit; can be interpreted independently of the other units of the text.
			\item \optword{Satellite (S)}: dependent unit; can only be interpreted using the nucleus.
		\end{itemize}
		\item Actors
		\begin{itemize}
			\item \optword{Reader (R)}: who reads the script.
			\item \optword{Writer (Sc)}: who wrote the script.
		\end{itemize}
		\item Relation definition
		\begin{itemize}
			\item Constraints on the nucleus
			\item Constraints on the satellite 
			\item Constraints the combination of both 
			\item Produced effect on either the reader, the writer or both.
		\end{itemize}
	\end{itemize}
	
\end{frame}

\begin{frame}
	\frametitle{Discourse coherence: Coherence relations}
	\framesubtitle{RST: RST relations (Elaboration)}
	
	\begin{itemize}
		\item \optword{Elaboration}: S gives additional information about the situation presented in N.
		\begin{itemize}
			\item \textbf{Constraints on N + S}: S presents additional information regarding the situation or some aspect of the theme presented in N.
			\item \textbf{Effect}: R recognizes the situation presented in S as providing additional information about N. 
			\item \textbf{Locus of the effect}: N + S
			\item E.g. \expword{[\textsubscript{N} The exam is easy.] [\textsubscript{S} It only takes an hour.]}
		\end{itemize}
	\end{itemize}
	
\end{frame}

\begin{frame}
	\frametitle{Discourse coherence: Coherence relations}
	\framesubtitle{RST: RST relations (Evidence)}
	
	\begin{itemize}
		\item \optword{Evidence}: S gives additional information about the situation presented in N in order to convince R to accept the information of N
		\begin{itemize}
			\item \textbf{Constraints on N}: R may not believe in N to a degree satisfactory to W.
			\item \textbf{Constraints on S}: R believes S or will find it credible
			\item \textbf{Constraints on N + S}: R's comprehending S increases R's belief of N
			\item \textbf{Effect}: R's belief of N is in increased
			\item \textbf{Locus of the effect}: N
			\item E.g. \expword{[\textsubscript{N} Kevin must be here.] [\textsubscript{S} His car is parked outside.]}
		\end{itemize}
	\end{itemize}
	
\end{frame}

\begin{frame}
	\frametitle{Discourse coherence: Coherence relations}
	\framesubtitle{RST: RST relations (Contrast)}
	
	\begin{itemize}
		\item \optword{Contrast}: An opposition relation between two nuclei.
		\begin{itemize}
			\item \textbf{Constraints on N}: multi-nuclear
			\item \textbf{Constraints on N + N}: no more than two nuclei; the situations presented in those two nuclei are (a) comprehended as the same in many aspects (b) comprehended as differing in a few aspects and (c) compared with respect to one or more of these differences.
			\item \textbf{Effect}: R recognizes the comparability and the difference(s) yielded by the comparison being made.
			\item \textbf{Locus of the effect}: multiple nuclei.
			\item E.g. \expword{[\textsubscript{N} He voted "No" to the new constitution.] [\textsubscript{N} His brother voted "Yes".]}
		\end{itemize}
	\end{itemize}
	
\end{frame}

\begin{frame}
	\frametitle{Discourse coherence: Coherence relations}
	\framesubtitle{RST: RST relations (Antithesis)}
	
	\begin{itemize}
		\item \optword{Antithesis}: An opposition relation between N and S so that R has a positive attitude towards N.
		\begin{itemize}
			\item \textbf{Constraints on N}: W has positive regard for the situation presented in N.
			\item \textbf{Constraints on N + N}: the situations presented in N and S are in contrast (cf. CONTRAST).
			\item \textbf{Effect}: R's positive regard for N is increased.
			\item \textbf{Locus of the effect}: N
			\item Ex. \expword{[\textsubscript{N} I defended the republic when I was young;] [\textsubscript{S} I won't give it up now that I'm old.]}
		\end{itemize}
	\end{itemize}
	
\end{frame}

\begin{frame}
	\frametitle{Discourse coherence: Coherence relations}
	\framesubtitle{RST: Subject matter vs. presentational relations}
	
	\begin{center}
		\footnotesize
	\begin{tblr}{
			colspec = {p{.4\textwidth}lp{.4\textwidth}},
			row{odd} = {lightblue},
			row{even} = {lightyellow},
			row{1} = {darkblue},
			rowsep = 0pt,
		}

		\bfseries\textcolor{white}{Subject matter} && \bfseries\textcolor{white}{Presentational}\\
		
		Elaboration
		
		Circumstance
		
		Solutionhood
		
		Volitional Cause
		
		Volitional Result
		
		Non-Volitional Cause
		
		Non-Volitional Result
		
		Purpose
		
		Condition
		
		Otherwise
		
		Interpretation
		
		Evaluation
		
		Restatement
		
		Summary
		
		Sequence
		
		Contrast
		
		&&
		Motivation
		
		Antithesis
		
		Background
		
		Enablement
		
		Evidence
		
		Justify
		
		Concession\\
	\end{tblr}
	\end{center}
	
\end{frame}

\begin{frame}
	\frametitle{Discourse coherence: Coherence relations}
	\framesubtitle{RST: semantic vs. pragmatic vs. textual relations}
	
	
	\begin{center}
		\footnotesize
	\begin{tblr}{
			colspec = {p{.3\textwidth}lp{.25\textwidth}lp{.2\textwidth}},
			row{odd} = {lightblue},
			row{even} = {lightyellow},
			row{1} = {darkblue},
			rowsep = 0pt,
		}

		\bfseries\textcolor{white}{Semantic} && \bfseries\textcolor{white}{Pragmatic} && \bfseries\textcolor{white}{Textual} \\
		
		Elaboration
		
		Circumstance
		
		Solutionhood
		
		Volitional Cause
		
		Volitional Result
		
		Non-Volitional Cause
		
		Non-Volitional Result
		
		Purpose
		
		Condition
		
		Interpretation
		
		Evaluation
		
		Sequence
		
		Contrast
		
		&&
		
		Motivation
		
		Antithesis
		
		Facilitation
		
		Evidence 
		
		Justification
		
		Concession
		
		&&
		
		Background
		
		Reformulation
		
		Summary\\
	\end{tblr}
	\end{center}
	
\end{frame}

\subsection{Penn Discourse TreeBank (PDTB)}

\begin{frame}
	\frametitle{Discourse coherence: Coherence relations}
	\framesubtitle{Penn Discourse TreeBank (PDTB)}
	
	\begin{itemize}
		\item Annotation of discourse connectives
		\begin{itemize}
			\item \optword{Subordination conjunctions}:  
			temporal (E.g., \expword{'when', 'as soon as'}), 
			causal (E.g., \expword{'because'}), 
			concessive (E.g., \expword{'although', 'even though'}), 
			objective (E.g.,\expword{'so that', 'in order that'}) et 
			conditional (E.g., \expword{'if', 'unless'}).
			
			\item \optword{Coordination conjunctions}: \expword{'and', 'but', 'or'}
			
			\item \optword{Adverbial conjunctions}: adverbs that express a discourse relation between events or states. E.g., \expword{'however', 'therefore', 'then', etc.}
			Prepositional phrases are also included in this class. E.g. \expword{'as a result',
			'in addition', 'in fact', etc. }
			
			\item \optword{Implicit conjunctions}: identified between two adjacent sentences that are not connected by explicit connectives.
		\end{itemize}
		\item Arguments annotation
		\begin{itemize}
			\item ARG2: this is the clause where the connective is bound (explicit case)
			\item ARG1: the other clause
		\end{itemize}
	\end{itemize}
	
\end{frame}

\begin{frame}
	\frametitle{Discourse coherence: Coherence relations}
	\framesubtitle{Penn Discourse TreeBank (PDTB): Connectives senses}
	
	\begin{table}
	\tiny\bfseries
	\begin{tblr}{
			colspec = {p{.1\textwidth}lp{.8\textwidth}},
			row{odd} = {lightblue},
			row{even} = {lightyellow},
			row{1} = {darkblue},
			rowsep = 1pt,
		}

		\bfseries\textcolor{white}{Connective} && \bfseries\textcolor{white}{Sens}\\
		
		after && succession (523), succession-reason (50), other (4) \\
		since && reason (94), succession (78), succession-reason (10), other (2) \\
		when && Synchrony (477), succession (157), general (100), succession-reason (65), Synchrony-general (50),
		Synchrony-reason (39), hypothetical (11), implicit assertion (11), Synchrony-hypothetical (10), other
		(69) \\
		while && juxtaposition (182), Synchrony (154), Contrast (120), expectation (79), opposition (78), Conjunction
		(39), Synchrony-juxtaposition (26), Synchrony-Conjunction (21), Synchrony-Contrast(22), COMPARISON (18), Synchrony-opposition (11), other (31) \\
		meanwhile && Synchrony-Conjunction (92), Synchrony (26), Conjunction (25), Synchrony-juxtaposition (15),
		other(35)\\
		but && Contrast (1609), juxtaposition (636), contra-expectation (494), COMPARISTON (260), opposition
		(174), Conjunction (63), Conjunction-Pragmatic contrast (14), Pragmatic-contrast (14), other (32)
		however Contrast (254), juxtaposition (89), contra-expectation (70), COMPARISON (49), opposition (31),
		other (12)\\
		although && expectation (132), Contrast (114) juxtaposition (34), contra-expectation (21), COMPARISON (16),
		opposition (9), other (2)\\
		and && Conjunction (2543), List (210), result-Conjunction (138), result (38), precedence-Conjunction (30),
		juxtaposition (11), other(30)\\
		if && hypothetical (682), general (175), unreal present (122), factual present (73), unreal past (53), expectation (34), implicit assertion (29), relevance (20), other (31)\\
	\end{tblr}
	\caption{Some connectives and statistics on their senses \cite{2008-prasad-al}}
	\end{table}
	
\end{frame}

\begin{frame}
	\frametitle{Discourse coherence: Coherence relations}
	\framesubtitle{Penn Discourse TreeBank (PDTB): Connectives categories}
	
	\begin{figure}
		\centering\vspace{-6pt}
		\begin{tcolorbox}[colback=white, colframe=blue, boxrule=1pt, text width=.9\textwidth]
		\begin{minipage}{.3\textwidth}
			\tiny\bfseries\vspace{-6pt}
			\begin{itemize}
				\item CONTINGENCY
				\begin{itemize}\tiny\bfseries
					\item Cause
					\begin{itemize}\tiny\bfseries
						\item Reason
						\item Result
					\end{itemize}
					\item Pragmatic Cause
					\begin{itemize}\tiny\bfseries
						\item Justification
					\end{itemize}
					\item Condition
					\begin{itemize}\tiny\bfseries
						\item Hypothetical
						\item General
						\item Unreal Present
						\item Unreal Past
						\item Factual Present
						\item Factual Past
					\end{itemize}
					\item Pragmatic Condition
					\begin{itemize}\tiny\bfseries
						\item Relevance
						\item Implicit Assertion
					\end{itemize}
				\end{itemize}
			\end{itemize}
		\end{minipage}
		\begin{minipage}{.3\textwidth}
			\tiny\bfseries
			\begin{itemize}\tiny\bfseries
				\item TEMPORAL
				\begin{itemize}\tiny\bfseries
					\item Asynchronous
					\item Synchronous
					\begin{itemize}\tiny\bfseries
						\item Precedence
						\item Succession
					\end{itemize}
				\end{itemize}
				\item COMPARISON
				\begin{itemize}\tiny\bfseries
					\item Contrast
					\begin{itemize}\tiny\bfseries
						\item Juxtaposition 
						\item Opposition
					\end{itemize}
					\item Pragmatic Contrast
					\item Concession
					\begin{itemize}\tiny\bfseries
						\item Expectation
						\item Contra-expectation
					\end{itemize}
					\item Pragmatic Concession
				\end{itemize}
			\end{itemize}
		\end{minipage}
		\begin{minipage}{.3\textwidth}
			\tiny\bfseries
			\begin{itemize}\tiny\bfseries
				\item EXPANSION
				\begin{itemize}\tiny\bfseries
					\item Conjunction
					\item Instantiation
					\item Restatement
					\begin{itemize}\tiny\bfseries
						\item Specification
						\item Equivalence
						\item Generalization
					\end{itemize}
					\item Alternative
					\begin{itemize}\tiny\bfseries
						\item Conjunction
						\item Disjunction
						\item Chosen Alternative
					\end{itemize}
					\item Exception
					\item List
				\end{itemize}
			\end{itemize}
		\end{minipage}\vspace{-6pt}
		\end{tcolorbox}\vspace{-6pt}
		\caption{PDTB sens hierarchy \cite{2008-prasad-al}}
	\end{figure}
	
\end{frame}

%===================================================================================
\section{Discourse Structure Parsing}
%===================================================================================

\begin{frame}
	\frametitle{Discourse coherence}
	\framesubtitle{Discourse Structure Parsing}
	
	\begin{itemize}
		\item A discourse is coherent if there are relations between its parts.
		\item The structure:
		\begin{itemize}
			\item Relations tree: RST
			\item Binary relations between each two parts: PDTB
		\end{itemize}
	\end{itemize}
	
\end{frame}

\subsection{RST parsing}

\begin{frame}
	\frametitle{Discourse coherence: Discourse Structure Parsing}
	\framesubtitle{RST parsing}
	
	\begin{itemize}
		\item Build a discourse relations tree
		\item Two steps: 
		\begin{enumerate}
			\item Elementary discourse units (EDUs) detection
			\item Relations classification
		\end{enumerate}
	\end{itemize}

	\begin{center}
		\hgraphpage[0.5\textwidth]{RST-tree.pdf}
	\end{center}
	
\end{frame}

\begin{frame}
	\frametitle{Discourse coherence: Discourse Structure Parsing}
	\framesubtitle{RST parsing: Elementary discourse units (EDUs) detection}
	
	\begin{itemize}
		\item Elementary discourse units (EDUs)
		\item A sentence can have many EDUs 
		\item E.g. \expword{[Mr. Rambo says]\textsubscript{e1} [that a 3.2-acre property]\textsubscript{e2} [overlooking the San Fernando Valley]\textsubscript{e3} [is priced at \$4 million]\textsubscript{e4} [because the late actor Erroll Flynn once lived there.]\textsubscript{e5}}
		\item Methods 
		\begin{itemize}
			\item \optword{Rule-based}
			\begin{itemize}
				\item parse tree
				\item surface features: punctuation and lexical markers such as connectives
			\end{itemize}
			\item Statistical
			\begin{itemize}
				\item \optword{Features-based}: using syntactic information and surface features to learn how to limit UEDs
				\item \optword{Embeddings-based}: text as sequence. E.g., \url{https://github.com/PKU-TANGENT/NeuralEDUSeg}
			\end{itemize}
		\end{itemize}
	\end{itemize}
	
\end{frame}

\begin{frame}
	\frametitle{Discourse coherence: Discourse Structure Parsing}
	\framesubtitle{RST parsing: EDUs detection by embeddings}
	
	\begin{figure}
		\centering
		\hgraphpage[.45\textwidth]{EDU_seg.pdf}
		\caption{EDUs segmentation proposed by \cite{2018-wang-al}; figure inspired by \cite{2019-jurafsky-martin}}
	\end{figure}
	
\end{frame}

\begin{frame}
	\frametitle{Discourse coherence: Discourse Structure Parsing}
	\framesubtitle{RST parsing: Relations classification}
	
	\begin{minipage}{.6\textwidth}
	\begin{itemize}
		\item Most used method: SHIFT-REDUCE
		\item Configuration: $C = (\sigma, \beta, A)$
		\begin{itemize}
			\item $\sigma$ is a stack
			\item $\beta$ is an input buffer
			\item $A$ is the list of created arcs (relations)
			\item $C_{initiale} = (\varnothing, w, \emptyset)$
			\item $C_{finale} = (\varnothing, \varnothing, A)$
		\end{itemize}
	\end{itemize}
	\end{minipage}
	\begin{minipage}{.38\textwidth}
		\hgraphpage{RST-transitions.pdf}
	\end{minipage}
	\begin{itemize}
		\item Operations 
		\begin{itemize}
			\item \optword{Shift}: put the first element of the buffer on the stack.
			\item \optword{Reduce}\textbf{(l, d)}: merge the top two subtrees of the stack, where \textbf{l} is the label of coherence relation, and \textbf{d} is the direction of nuclearity: \textbf{d $ \in $ \{NN, NS, SN\}}.
			\item \optword{Pop Root}: delete the top tree from the stack.
		\end{itemize}
		\item Training 
		\begin{itemize}
			\item Features-based
			\item Embeddings-based
		\end{itemize}
	\end{itemize}
	
\end{frame}

\begin{frame}
	\frametitle{Discourse coherence: Discourse Structure Parsing}
	\framesubtitle{RST parsing: Relations classification (example)}
	
	\begin{figure}
		\begin{minipage}{0.25\textwidth}
			\hgraphpage{RST_exp.pdf}
		\end{minipage}
		\begin{minipage}{0.70\textwidth}
			\scriptsize
			\begin{itemize}
				\item e\textsubscript{1}: American Telephone \& Telegraph Co. said it
				\item e\textsubscript{2}: will lay off 75 to 85 technicians here , effective Nov. 1.
				\item e\textsubscript{3}: The workers install , maintain and repair its private branch exchanges,
				\item e\textsubscript{4}: which are large intracompany telephone networks.
			\end{itemize}
		\end{minipage}
		
		\small
		\begin{tabular}{lllll}
			\hline\hline 
			Step & Stack & Buffer & Action & Added Relation \\
			\hline
			01 & $\emptyset$ & $e_1, e_2, e_3, e_4$ & SH & / \\
			02 & $e_1$ & $e_2, e_3, e_4$ & SH & / \\
			03 & $e_1, e_2$ & $e_3, e_4$ & RD(attr, SN) & $\widehat{e_1 \mathbf{e_2}}$ \\
			04 & $e_{1:2}$ & $e_3, e_4$ & SH & / \\
			05 & $e_{1:2}, e_3$ & $e_4$ & SH & / \\
			06 & $e_{1:2}, e_3, e_4$ & $\emptyset$ & RD(elab, NS) & $\widehat{\mathbf{e_3} e_4}$ \\
			07 & $e_{1:2}, e_{3:4}$ & $\emptyset$ & RD(elab, SN) & $\widehat{e_{1:2}\mathbf{e_{3:4}}}$ \\
			08 & $e_{1:4}$ & $\emptyset$ & PR & / \\
			\hline\hline
		\end{tabular}
		\caption{Example of RST parsing using Shif-Reduce \cite{2018-yu-al}}
	\end{figure}
	
\end{frame}

\begin{frame}
	\frametitle{Discourse coherence: Discourse Structure Parsing}
	\framesubtitle{RST parsing: Relations classification using embeddings (1)}
	
	\begin{itemize}
		\item Case study: \cite{2018-yu-al} (Encoder-decoder)
		\item Code: \url{https://github.com/yunan4nlp/NNDisParser}
		\item Encoder 
		\begin{itemize}
			\item Encode words: $ w_i $ a word with a PoS $t_i$
			\[x_i^w = embedding(w_i) \oplus embedding(t_i)\]
			\item Represent sequential words
			\[ \{h_1^w, h_2^w, \ldots, h_m^w \} = biLSTM(\{x_1^w, x_2^w, \ldots, x_m^w \})\]
			\item Represent one EDU $\{w_s, w_{s+1}, \ldots, w_t \}$
			\[ x^e = \frac{1}{t-s+1} \sum_{k=s}^{t} h_k^w\]
			\item Represent sequential EDU
			\[ \{h_1^e, h_2^e, \ldots, h_n^e \} = biLSTM(\{x_1^e, x_2^e, \ldots, x_n^e \})\]
		\end{itemize}
	\end{itemize}
	
\end{frame}

\begin{frame}
	\frametitle{Discourse coherence: Discourse Structure Parsing}
	\framesubtitle{RST parsing: Relations classification using embeddings (2)}
	
	\begin{itemize}
		\item Decoder 
		\begin{itemize}
			\item A feed-forward neural layer to infer the action $o$
			\[o = W(h_{s0}^{sbt} \oplus h_{s1}^{sbt} \oplus h_{s2}^{sbt} \oplus h_{q0}^{e})\]
			\item $ h_{e}^{q0} $: sequential representation of the first EDU in the buffer
			\item $h_{si}^{sbt}$: sequential representation of the subtree $i$ in the stack
			\item the representation of a subtree between two EDUs $ s= \{e_i, \ldots, e_j\}$ is the average of the covered EDUs representations. 
			\[ h_{s}^{sbt} = \frac{1}{j-i+1} \sum_{k=i}^{j} h_k^e\]
		\end{itemize}
	\end{itemize}
	
\end{frame}

\subsection{PDTB parser}

\begin{frame}
	\frametitle{Discourse coherence: Discourse Structure Parsing}
	\framesubtitle{PDTB parser}
	
	\begin{itemize}
		\item Find the relation between two spans
		\item 4 steps: 
		\begin{enumerate}
			\item Find the discourse connectives
			\begin{itemize}
				\item Using sens disambiguation (delimiter or not)
			\end{itemize}
			\item Find the two spans for each connective
			\begin{itemize}
				\item Given a span with a connector, adjacent spans can be classified as having a relation or not (the same manner as EDUs classification)
			\end{itemize}
			\item Label the relationship between these spans (explicit connectives)
			\begin{itemize}
				\item Connectives are labels
			\end{itemize}
			\item Assign a relation between every adjacent pair of sentences (implicit connectives)
			\begin{itemize}
				\item For example, using BERT with the two spans as input, the output of $ <CLS> $ can be followed by a MLP to infer the relation
			\end{itemize}
		\end{enumerate}
	\end{itemize}
	
\end{frame}

\begin{frame}
	\frametitle{Discourse coherence}
	\framesubtitle{Discourse Structure Parsing: Some humor}
	
	\begin{center}
		\vgraphpage{humor/humor2.jpg}
	\end{center}
	
\end{frame}

%===================================================================================
\section{Entity-based coherence}
%===================================================================================

\begin{frame}
	\frametitle{Discourse coherence}
	\framesubtitle{Entity-based coherence}
	
	\begin{itemize}
		\item A discourse is coherent if it discusses a certain entity
		\item In all the discourse, this entity must remain the most important.
	\end{itemize}

	\begin{exampleblock}{Example of entity-based coherence \cite{2019-jurafsky-martin}}
		\small
		\begin{minipage}{.48\textwidth}
			\begin{enumerate}
				\item John went to his favorite music store to buy a piano. [John]
				\item He had frequented the store for many years. [John]
				\item He was excited that he could finally buy a piano. [John]
				\item He arrived just as the store was closing for the day. [John]
			\end{enumerate}
		\end{minipage}
		\begin{minipage}{.48\textwidth}
			\begin{enumerate}
				\item John went to his favorite music store to buy a piano. [John]
				\item It was a store John had frequented for many years. [The store]
				\item He was excited that he could finally buy a piano. [John]
				\item It was closing just as John arrived. [The store]
			\end{enumerate}
		\end{minipage}
	\end{exampleblock}
	
\end{frame}

\subsection{Centering theory}

\begin{frame}
	\frametitle{Discourse coherence: Entity-based coherence}
	\framesubtitle{Centering theory: Discourse center}
	
	\begin{itemize}
		\item $\mathbf{U_n}$: an utterance
		\item $\mathbf{C_b(U_n)}$: \optword{Backward-looking center}
		\begin{itemize}
			\item the current salient entity
			\item the one being focused on in the discourse after $ U_{n-1} $.
		\end{itemize}
		\item $\mathbf{C_f(U_n)}$ : \optword{Forward-looking center}
		\begin{itemize}
			\item potentially salient entities in the future
			\item those candidates to be $C_b(U_{n+1})$.
		\end{itemize}
		\item $\mathbf{C_p(U_n) \in C_f(U_n)}$ : \optword{Preferred center}
		\begin{itemize}
			\item candidate center to be $C_b(U_{n+1})$
			\item score based on grammatical role (subjects are more important than objects than the rest), order (E.g. \expword{In Arabic, what is first is more important}), etc.
		\end{itemize}
		\item $C_b(U_n)$ must be represented by a pronoun in $U_{n+1}$
	\end{itemize}
	
\end{frame}

\begin{frame}
	\frametitle{Discourse coherence: Entity-based coherence}
	\framesubtitle{Centering theory: Transitions}
	
	\begin{itemize}
		\item Transitions order from the most coherent
		\item \optword{CONTINUE}: the speaker talks about an entity and has intention to speak about it in the future
		\item \optword{RETAIN}: the speaker talks about an entity and has intention to change it in the future
		\item \optword{SHIFT}: the speaker changed the central entity
		\begin{itemize}
			\item \textbf{Smooth-SHIFT}: after shifting, he has the intention to talk about it in the future
			\item \textbf{Rough-SHIFT} : after shifting, he has the intention to change it in the future
		\end{itemize}
	\end{itemize}

	\begin{center}
		\tiny\bfseries
		\begin{tblr}{
				colspec = {p{.2\textwidth}p{.2\textwidth}p{.2\textwidth}},
				row{odd} = {lightblue},
				row{even} = {lightyellow},
				row{1} = {darkblue},
				rowsep = 2pt,
			}
			
			& \bfseries\color{white}$\mathbf{C_b(U_n) = C_b(U_{n-1})}$
			
			OU $\mathbf{C_b(U_n) = NULL}$
			& \bfseries\color{white}$\mathbf{C_b(U_n) \ne C_b(U_{n-1})}$\\
			
			$\mathbf{C_b(U_n) = C_p(U_n)}$ &
			CONTINUE & Smooth-SHIFT\\
			
			$\mathbf{C_b(U_n) \ne C_p(U_n)}$ &
			RETAIN & Rough-SHIFT\\
		\end{tblr}
	\end{center}
	
\end{frame}

\subsection{Entity Grid model}

\begin{frame}
	\frametitle{Discourse coherence: Entity-based coherence}
	\framesubtitle{Entity Grid model}
	
	\begin{itemize}
		\item A document is represented by a matrix
		\begin{itemize}
			\item Line: sentences
			\item Columns: entities
			\item Values: Grammatical function of the entity; subject (S), Object (O), Other (X) or the entity does nor exist (\_).
		\end{itemize}
		\item A document is coherent, if it follows a pattern
		\begin{itemize}
			\item This pattern can be represented by the frequencies of grammatical transitions. E.g. \expword{SS, SO, SX, S\_, OS, OO, OX, O\_, ...}
			\item Machine learning is used to judge the coherence of a document
		\end{itemize}
		
	\end{itemize}
	
\end{frame}

\begin{frame}
	\frametitle{Discourse coherence: Entity-based coherence}
	\framesubtitle{Entity Grid model: Example of document's representation}
	
	\begin{figure}
		\centering\tiny\bfseries
		\begin{minipage}{.8\textwidth}
			\begin{enumerate}
				\item\ [The Justice Department]\textsubscript{s} is conducting an [anti-trust trial]\textsubscript{o} against [Microsoft Corp.]\textsubscript{x} with [evidence]\textsubscript{x} that [the company]\textsubscript{s} is increasingly attempting to crush [competitors]\textsubscript{o}.
				\item\ [Microsoft]\textsubscript{o} is accused of trying to forcefully buy into [markets]\textsubscript{x} where [its own products]\textsubscript{s} are not competitive enough to unseat [established brands]\textsubscript{o}.
				\item\ [The case]\textsubscript{s} revolves around [evidence]\textsubscript{o} of [Microsoft]\textsubscript{s} aggressively pressuring [Netscape]\textsubscript{o} into merging [browser software]\textsubscript{o}.
				\item\ [Microsoft]\textsubscript{s} claims [its tactics]\textsubscript{s} are commonplace and good economically.
				\item\ [The government]\textsubscript{s} may file [a civil suit]\textsubscript{o} ruling that [conspiracy]\textsubscript{s} to curb [competition]\textsubscript{o} through [collision]\textsubscript{x} is [a violation of the Sherman Act]\textsubscript{o}.
				\item\ [Microsoft]\textsubscript{s} continues to show [inscreased earnings]\textsubscript{o} despite [the trial]\textsubscript{x}.
			\end{enumerate}
		\end{minipage}
		
		\color{blue}
		\begin{tabular}{ccccccccccccccccc}
			& \rotatebox[origin=c]{90}{Department} & \rotatebox[origin=c]{90}{Trial} & \rotatebox[origin=c]{90}{Microsoft} &
			\rotatebox[origin=c]{90}{Evidence} & \rotatebox[origin=c]{90}{Competitors} & \rotatebox[origin=c]{90}{Markets} &
			\rotatebox[origin=c]{90}{Products} & \rotatebox[origin=c]{90}{Brands} & \rotatebox[origin=c]{90}{Case} &
			\rotatebox[origin=c]{90}{Netscape} & \rotatebox[origin=c]{90}{Software} & \rotatebox[origin=c]{90}{Tactics} &
			\rotatebox[origin=c]{90}{Government} & \rotatebox[origin=c]{90}{Suit} & \rotatebox[origin=c]{90}{Earnings} & \\
			
			1. & s & o & s & x & o & - & - & - & - & - & - & - & - & - & - & 1. \\
			2. & - & - & o & - & - & x & s & o & - & - & - & - & - & - & - & 2. \\
			3. & - & - & s & o & - & - & - & - & s & o & o & - & - & - & - & 3. \\
			4. & - & - & s & - & - & - & - & - & - & - & - & s & - & - & - & 4. \\
			5. & - & - & - & - & - & - & - & - & - & - & - & - & s & o & - & 5. \\
			6. & - & x & s & - & - & - & - & - & - & - & - & - & - & - & o & 6. \\
			
		\end{tabular}
		\caption{Example of document representation using entities \cite{2008-barzilay-lapata}}
	\end{figure}
	
\end{frame}

\begin{frame}
	\frametitle{Discourse coherence: Entity-based coherence}
	\framesubtitle{Entity Grid model: Transitions representation}
	
	\begin{itemize}
		\item A document is represented by grammatical transitions probabilities
		\item A probability is the ratio between the number of a transition and the number of all transitions
		\item E.g. In the past example: \expword{P(S\_) = 2/75}
	\end{itemize}
	\begin{figure}
		\tiny\bfseries
		\begin{tabular}{lllllllllllllllll}
			\hline
			& s s & s o & s x & s - & o s & o o & o x & o - & x s & x o & x x & x - & - s & - o & - x & - - \\
			\hline
			d1 & .01 & .01 & .00 & .08 & .01 & .00 & .00 & .09 & .00 & .00 & .00 & .03 & .05 & .07 & .03 & .59 \\
			d2 & .02 & .01 & .01 & .02 & .00 & .07 & .00 & .02 & .14 & .14 & .06 & .04 & .03 & .07 & .01 & .36 \\
			d3 & .02 & .00 & .00 & .03 & .09 & .00 & .09 & .06 & .00 & .00 & .00 & .05 & .03 & .07 & .17 & .39 \\
			\hline
		\end{tabular}
		\caption{Example of a documents representation by grammatical transitions \cite{2008-barzilay-lapata}}
	\end{figure}
	
\end{frame}

\begin{frame}
	\frametitle{Discourse coherence: Entity-based coherence}
	\framesubtitle{Entity Grid model: Training}

	\begin{itemize}
		\item The algorithm learns to rate a text according to its coherence
		\item Incoherent texts must have lower scores than coherent ones
		\item Using coherence-annotated texts by experts
		\item Using a self-supervision method: create inconsistent text from consistent text
		\begin{itemize}
			\item Sentence order discrimination: we compare a document to a random permutation of its sentence. 
			\item Sentence insertion: we modify the order of a single sentence to have many incoherent texts.
			\item Sentence order reconstruction: we take a document, randomize the sentences, and train the model to put them back in the correct order. 
		\end{itemize}
	\end{itemize}
	
\end{frame}

\begin{frame}
	\frametitle{Discourse coherence}
	\framesubtitle{Entity-based coherence: Some humor}
	
	\begin{center}
		\vgraphpage{humor/humor3.png}
	\end{center}
	
\end{frame}


\insertbibliography{NLP10}{*}

\end{document}

